\label{sec:implementation}

The implementation assigns local differentiable kernels to JAX and distributed
algebraic objects to PETSc. Each MPI rank stores one contiguous interval of the
free coefficient vector and the ghost values required by its overlap elements.
The same ownership map is used for the residual, tangent rows, Riesz operator,
and canonical state export. This common map is essential: comparing local
arrays before undoing a route-specific permutation would not compare the same
algebraic object.

\subsection{Local derivative routes}
\label{subsec:implementation-routes}

The element route evaluates a scalar $\Phi_e$ and obtains its gradient and
Hessian by JAX transformations. The constitutive route evaluates the same
selected quadrature density, obtains $\sigma_{eq}$ and $C_{eq}$ from
\eqref{eq:constitutive-objects}, and contracts them according to
\eqref{eq:branchwise-route-identity}. For the endpoint-surrogate experiments, the
residual is also evaluated by element AD; only the tangent differentiation
boundary changes. This choice prevents a residual-formula difference from
being hidden inside a tangent comparison.

The colored route first builds the complete structural pattern of every owned
row. For the three-dimensional vector problem, it colors the active scalar-node
support at distance two with deterministic lexical weights and lifts a scalar
color $c_s(a)$ to displacement component $k$ as $3c_s(a)+k$. The lift is
validated against every scalar support row before differentiation. It is a
conservative coloring of the vector column-interference graph by
Proposition~\ref{prop:component-lifted-coloring}, including under componentwise
constraints. Exact JAX Hessian-vector products are then evaluated in fixed
batches of four colors, and each batch is inserted into the uniquely identified
owned entries before the next batch is formed. The final incomplete batch is
zero padded, so compilation shape does not depend on the color count. The
implementation does not assume that different ranks obtain the same color
labels. It relies instead on Proposition~\ref{prop:owned-row-colored-recovery}
and unique PETSc row ownership.

All three routes use binary64 arithmetic. Compilation and kernel warmup occur
before measured repetitions. Compiled element, constitutive, and action kernels
remain separate because their traces and intermediate arrays are different
scientific costs. At the feasible one-rank $P_1$ route points, full-mesh
compressed sparse row (CSR) arrays are written to transient disk-backed storage
and compared in bounded chunks; the complete matrix norms and maximum entrywise
defects are evaluated rather than sampled. Other full-mesh route points use four
prescribed tangent actions, while the fixed-element verifiers compare complete
local Hessians.

\subsection{Rank-local mesh and sparse assembly}
\label{subsec:rank-local-assembly}

Let $I_r$ denote rank $r$'s owned free DOFs and $G_r$ its required ghosts. The
retained overlap is
\begin{equation}
  \mathcal{T}_r^{\mathrm{loc}}
  :=\{e:\operatorname{dofs}(e)\cap I_r\ne\varnothing\}.
  \label{eq:rank-overlap-elements}
\end{equation}
Before each local evaluation, a point-to-point exchange populates $G_r$ and
the affine lift supplies constrained values. Residual and tangent contributions
are accumulated only for rows in $I_r$; off-rank columns retain global indices.
PETSc's distributed coordinate format is preallocated once and reused.

Scalar energy requires a separate ownership rule because overlap elements occur
on more than one rank. Each element has a deterministic owner $o(e)$ and weight
$\omega_{r,e}=1$ if $o(e)=r$. Hence
\begin{equation}
  \sum_r\sum_{e\in\mathcal{T}_r^{\mathrm{loc}}}
  \omega_{r,e}\Phi_e
  =\sum_{e\in\mathcal{T}_h}\Phi_e.
  \label{eq:owned-element-energy}
\end{equation}
This identity, the affine lift, the canonical free-DOF order, and the CSR
indices are included in the distribution-equivalence test.

The distributed hyperelastic element path builds the structured mesh and overlap
data rank locally. Its one-rank replicated construction is retained only as an
algebraic comparator. The endpoint-surrogate meshes and material data are immutable
inputs; quadrature rules have stable names and are recorded in every result.

\subsection{Nonlinear and linear solvers}
\label{subsec:solver-implementation}

The nonlinear driver accepts either line-search Newton or a trust-region
policy \citep{nocedal2006numerical,conn2000trust}. Armijo backtracking is used for the controlled Newton comparisons. The
reduced-subspace trust method minimizes the quadratic model on the span of the
Newton and gradient directions and applies the same configured Armijo rule when
a post-step search is requested. Thus a method label fixes the actual
globalization algorithm rather than selecting an undocumented fallback.

Each accepted Newton direction is produced by a PETSc Krylov solve on the
current tangent. The result record contains the requested and effective
Krylov-solver tolerances, iteration count, convergence reason, and an
independently recomputed true residual. A nonconverged Krylov solve, nonfinite
direction, rejected line search, trust-region rejection cap, or nonlinear
iteration cap remains a failure. Such a row is never converted into a
successful observation with a large time.

The primary verification configurations use ordinary Krylov methods with
hypre preconditioning or a sparse direct factorization. No preconditioner
ranking is inferred in this paper.
Reference-elastic Riesz operators are independent of the Newton tangent and
remain fixed over the associated nonlinear trajectory.

\subsection{Riesz operators and terminal recomputation}
\label{subsec:riesz-implementation}

For a lumped scalar mass operator, the diagonal weights are assembled on the
exact free space and must be finite and strictly positive. For mechanics,
$K_{\mathrm{ref}}$ is copied from the exact discrete elastic tangent at the
declared reference state. Hyperelasticity uses the identity deformation
$y(X)=X$ with both end faces constrained. The branch-structured mechanics
Riesz operator uses the
elastic constitutive branch at zero displacement with the glued-bottom
constraint.

The constrained matrix is first checked for scale-aware symmetry. A sparse
direct factorization then reports its numerical inertia under the stored
factorization and scaling options; admission requires
\begin{equation}
  (n_-,n_0,n_+)=(0,0,n_{\mathrm{free}}).
  \label{eq:riesz-inertia-gate}
\end{equation}
The dual norm is evaluated by an independent conjugate-gradient (CG) solve with
a recomputed true-residual threshold. PETSc options cannot silently modify this
solve unless an explicit, namespaced override is recorded. At every terminal state, the driver
re-evaluates the energy, gradient, primal state norm, dual residual, and
correction from the final stored vector. A failed run therefore cannot inherit
an initial or pre-step residual as its endpoint value.

\subsection{Canonical records}
\label{subsec:canonical-records}

Every runner writes strict JavaScript Object Notation (JSON): nonfinite
sentinels are serialized as
\code{null}, and the file becomes visible only after an atomic replacement.
The schema records the experiment, case, route, repetition, MPI size, threads,
input and code identities, solver controls, terminal status, norms, work
counters, timing phases, and state hashes. Canonical NumPy \code{.npz} states
use the original global ordering and include coordinate, connectivity,
constraint, and metric metadata. State identities use the 256-bit Secure Hash
Algorithm (SHA-256) applied to the array data type, shape, and raw binary64
bytes. Hash equality is an exact reproducibility check, not a numerical norm or
an assertion of physical endpoint equivalence.
